\section{Thiết kế thực nghiệm}
\label{sec:experiment}

\subsection{Câu hỏi đánh giá và bộ màn chơi}

Thực nghiệm so sánh bốn thuật toán theo thời gian, bộ nhớ, số nút mở rộng và
chất lượng lời giải. Nhóm xem heuristic có giúp A* mở rộng ít nút hơn UCS và BFS
hay không, và chi phí không đều có khiến UCS chọn đường khác BFS hay không.
Trường hợp DFS được xem xét theo sự đánh đổi giữa tốc độ và chất lượng lời giải.
Các nhận xét chỉ dựa trên 15 màn chơi cố định, được chia theo cơ chế ở
Bảng~\ref{tab:level-groups}.

\begin{table}[H]
\centering
\caption{Các nhóm màn chơi dùng trong thực nghiệm}
\label{tab:level-groups}
\small
\begin{tabularx}{\textwidth}{@{}l l X@{}}
\toprule
\textbf{Nhóm} & \textbf{Màn chơi} & \textbf{Đặc trưng} \\
\midrule
Cơ bản & 01--04 & Floor, Void, Goal và ba tư thế khối \\
Nâng cao, không tách khối & 05, 06, 11, 12, 13 & Fragile, cầu, Soft/Heavy Switch \\
Tách và ghép khối & 07, 08, 14 & Tách khối, đổi khối đơn vị, ghép khối và cầu \\
Kết hợp & 09, 10, 15 & Fragile, công tắc, cầu và tách/ghép khối \\
\bottomrule
\end{tabularx}
\end{table}

\subsection{Môi trường thực nghiệm}

Thực nghiệm được tiến hành ngày \benchmarkdate\ trên máy tính chạy Linux
x86\_64, sử dụng bộ xử lý Intel Core i7-12650H và 15 GiB RAM. Chương trình
chạy bằng \benchmarkpython\ trong môi trường ảo Python. Toàn bộ thuật toán
được đánh giá trên cùng máy, cùng phiên bản mã nguồn tại commit
\texttt{7284a54} và cùng bộ 15 màn chơi.
Cấu hình phần cứng và phần mềm được giữ như nhau giữa các lượt đo nhằm giảm
sai khác do môi trường khi so sánh.

\subsection{Quy trình và kiểm soát công bằng}

Công cụ chạy thực nghiệm duyệt 15 màn chơi, bốn thuật toán BFS/DFS/UCS/A* và ba
lần lặp cho từng cặp, tổng cộng 180 lượt. Mỗi lượt tạo một \texttt{GameEngine}
mới và lấy trạng thái đầu. Các lượt được đo tuần tự để mỗi phiên tìm kiếm dùng
riêng \texttt{tracemalloc}. Không đặt giới hạn thời gian hoặc số nút. Các bộ
giải dùng cùng màn chơi, bộ luật, thứ tự sinh trạng thái kế tiếp và cách đo;
UCS/A* dùng chung hàm chi phí chuẩn. Thời gian dựng hình và phát lại không nằm
trong phép đo. CSV từng lượt lưu kết quả của mỗi lượt; CSV tổng hợp tính trung
bình và độ lệch chuẩn tổng thể bằng \texttt{statistics.pstdev} trên các lượt
\texttt{FOUND}.

\subsection{Chỉ số đánh giá}

Bốn nhóm chỉ số trong Bảng~\ref{tab:metrics} tách hiệu năng tìm kiếm khỏi chất
lượng đường đi. Nhóm chất lượng gồm ba đại lượng: số bước di chuyển, số hành
động và tổng chi phí.

\begin{table}[H]
\centering
\caption{Bốn nhóm chỉ số đánh giá thuật toán}
\label{tab:metrics}
\small
\renewcommand{\arraystretch}{1.18}
\begin{tabularx}{\textwidth}{@{}l>{\raggedright\arraybackslash}X@{}}
\toprule
\textbf{Chỉ số} & \textbf{Cách đo} \\
\midrule
Thời gian tìm kiếm & Đo bằng \texttt{perf\_counter}, đơn vị mili giây; không gồm tải màn chơi, dựng hình, ghi CSV hoặc phát lại \\
\addlinespace[2pt]
Bộ nhớ đỉnh & Lượng cấp phát Python cao nhất trừ mức ban đầu, đo bằng \texttt{tracemalloc}; không phải tổng bộ nhớ tiến trình (RSS) \\
\addlinespace[2pt]
Số nút mở rộng & Phần tử hợp lệ được lấy khỏi tập chờ xét và được dùng để sinh trạng thái kế tiếp; không tính trạng thái đích và phần tử lỗi thời \\
\addlinespace[2pt]
Chất lượng lời giải & \texttt{solution\_length} đếm hành động di chuyển; \texttt{action\_count} tính cả \textsc{SwitchCube}; \texttt{solution\_cost} dùng hàm chi phí của thuật toán \\
\bottomrule
\end{tabularx}
\end{table}

\subsection{Phương pháp xác minh lời giải}

Bộ kiểm thử tự động đạt 203/203 trường hợp trong môi
trường chạy này. Cả 180 lượt đo đều trả \texttt{FOUND}. Công cụ chạy thực
nghiệm chỉ ghi chỉ số, không lưu chuỗi hành động hoặc tự phát lại trong lượt
đo. Tập lệnh kiểm tra báo cáo đối chiếu 180 dòng CSV từng lượt, 60 dòng CSV
tổng hợp, ba lần lặp cho mỗi cặp và các giá trị trung bình, độ lệch chuẩn. Sau
đó, tập lệnh giải lại một lần cho từng cặp màn chơi--thuật toán và dùng
\texttt{verify\_solution} áp từng hành động qua \texttt{GameEngine} tới Goal:
60/60 lời giải hợp lệ. Thời gian xác minh độc lập không được cộng vào thời gian
tìm kiếm chính thức; luồng GUI xác minh trước khi phát lại như đã trình bày ở
Mục~\ref{sec:implementation}.
